\documentclass[11pt]{article}
\usepackage[T1]{fontenc}\usepackage{amsmath,amssymb,amsthm,lmodern,geometry}
\geometry{margin=.92in}\setlength{\emergencystretch}{2em}
\newtheorem{theorem}{Theorem}\newtheorem{lemma}{Lemma}
\newcommand{\dd}{\,d}\newcommand{\E}{\mathbb E}
\title{The critical right-half entry law and a resident-phase escape test}
\author{Standalone analytical increment}\date{Version 1, 3 November 2026}
\begin{document}\maketitle
\begin{abstract}
Positive-noise Q4 comparison ancestry supplies an aligned atom of mass
\(A/8\) at the resident. The complementary Q1 comparison law remains
distributed. For a prescribed compact leading entry we derive exact
moment equations and a self-consistent, finite-time escape criterion.
It gives an explicit positive mass of newly inactive labels and an
exponential lower barrier up to a specified radius. It applies, in
particular, to a saturated family of truncated critical profiles with
sufficiently small scale parameter. We identify the face inequality
that prevents this conclusion for unrestricted entry parameters.
\end{abstract}

\section{Domains}\label{an04rpsv1:domains}
Set \(\lambda=\rho^2,\ 13/20\le\rho\le7/10,\ \omega=(1-\lambda)/2,\ 
\ell=\log(1/q),\ S_0=cq^{1/\omega},\ T_A=\ell/\omega+\log(A/c)\).
Write \(a=\rho K(\rho a)>0,\ K(t)=\phi(t)+t\Phi(t)\),
\(d=[1-\lambda\Phi(\rho a)]/2\). All entry limits take q to zero
at fixed positive \(A,c\), uniformly also on compact positive ranges.
Offsets lie in a compact set separated above \(a\).
\begin{itemize}
\item PROFILE angle and coordinate concern the aligned Gaussian
target-only comparison with any positive initial radial weight. Its
local root estimate uses fixed bounded scaled charts at small q.
Its full Q4 corollary has the all-right-half time guard
\(T\ge(2/\omega)\ell+C,\ q^2T\le1\); the critical clock violates
the first condition. We derive the needed Q4 estimate directly.
\item CFS laws concerns Q1 comparison ancestry at \(T_A\), compact
interiors for individual limits, endpoint mass bounds for its integrated
law. CFS spread uses compact co-ordered labels in a single leading
field with \(N=0,\ 0<M\le1\). CFS sat starts below strong half mass.
\item WIT popsystem and P leading/order supply the defined leading
characteristics, with finite masses and bounded support on the time
interval in use. Comparison of two receivers holds with the same
time-dependent donors. We take \(N=0,\ 0<M_0\le1\) throughout the
resident calculation. At or near saturation the mass used is stated
explicitly. P spectra concerns infinitesimal offsets about a coherent
resident, not finite-amplitude escape. No WIT rate transfer or passage
from entry to a trained population is assumed.
\end{itemize}

\section{Q4 and the right-half comparison law}\label{an04rpsv1:entry}
\begin{theorem}\label{an04rpsv1:q4}
On compact interior Q4 labels \(-\pi/2<\beta<0\),
\[
 a-\widehat\theta(T_A,\beta)/q\asymp
 q^{d(1/\omega-2)},\qquad
 \widehat m(T_A,\beta)=A\cos^2\beta\,(1+o(1)).
\]
The Q4 radial mass is \(A/8+O(Aq^{1/4})\). Outside an endpoint
strip of mass \(O(Aq^{1/4})\), the absolute scaled displacement from
\(a\) is \(O(q^{d(1/\omega-5/2)}+q^2)\).
\end{theorem}
\begin{proof}
For \(r=-\theta>0\) the exact velocity from PROFILE angle is
\[
 2\theta'=\sin r\cos r+q\sin r K(-\cos r/q)
                      +\rho q\cos r K(-\rho\sin r/q).
\]
In particular \(r\) decreases. For an interior compact, until the
first hit of \(-qL\) at fixed large \(L\), put \(v=\tan r/q\).
Then \(-v'=v/2+E\) where, with \(D=\sqrt{1+q^2v^2}\),
\(2E=qvD K(-1/(qD))+\rho D K(-\rho v/D)\).
Here \(D\) is uniformly bounded, \(E\ge0\) and \(E\le C\).
Integrating the reciprocal speed gives \(2\ell+O(1)\):
subtract \(2/v\) on \(v\ge\tan(qL)/q\); the absolute difference is
at most \(4C/v^2\). No approximation with bounded \(w=\theta/q\)
is used at the initial angle of order one.

On \([-L,L]\), \(f_q(x)=\theta'(qx)/q
=[\rho K(\rho x)-x]/2+O_{C^2}(q^2)\) by the own-gate
Gaussian tail and Taylor expansion. The root \(a_q=a+O(q^2)\)
attracts from both sides and \(-f_q'(a_q)=d+O(q^2)\).
For \(h=a_q-x>0\), \(h'=-b_q(h)h\), where
\(b_q(h)\ge\omega/2\) and
\(|b_q(h)-b_q(0)|\le C h\). Integration first bounds
\(\int h\le2h(0)/\omega\), then bounds \(h(t)/h(0)\) above and below
by fixed multiples of \(\exp[-(d+O(q^2))t]\).
The time from \(-qL\) to observation is
\((1/\omega-2)\ell+O(1)\). Since \(0<d(1/\omega-2)<1\),
the \(O(q^2)\) shift cannot reverse the compact-interior sign.
In fact the entire quadrant stays below \(a\) at observation for small
q: the label initially at zero is the largest by scalar order and its
gap to \(a_q\) is comparable to \(\exp[-d\ell/\omega]\), with
\(d/\omega<2\). The label at the negative vertical axis is included
by continuity; no uniform lower arrival-speed assertion is needed there.

For a quantitative strip estimate restrict to \(\cos\beta\ge q^{1/4}\).
The first velocity inequality gives time to \(-qL\) at most
\(2\log[\tan|\beta|/\tan(qL)]\le(5/2)\ell+C\) whenever a hit is
needed. Labels starting inside the local chart have at least as much
relaxation time. This proves the asserted upper displacement, since
\(1/\omega>5/2\) uniformly.
With \(U_1=\sqrt m\cos\theta\), PROFILE coordinate gives exactly
\[
 (\log U_1)'=\tfrac12+d_q+
 \frac{qK(-\cos\theta/q)}{2\cos\theta},\qquad 0\le d_q\le q^2 .
\]
On this retained strip complement \(\cos\theta\ge q^{1/4}/2\)
through observation, by monotonicity towards the root. The final
cosine is \(1+O(q^2)\). Hence the relative mass error is
\(O(q^2(1+\ell))\). On either endpoint strip of length \(\eta\)
the gated covariance bound \(m'\le(1+2q^2)m\) costs \(O(A\eta)\).
The omitted negative-vertical strip has length \(O(q^{1/4})\);
the horizontal endpoint is handled by the local estimate. The
individual endpoints have zero mass. Integration against
\(d\beta/(2\pi)\) proves the mass formula.
\end{proof}

Put \(b=(c/A)^\omega\) and use CFS's monotone quadrature
\[
 B(w)=[w-\rho K(\rho w)]/2,\quad
 H(w)=w\exp\left[-\omega\int_w^\infty
       \left(B(r)^{-1}-(\omega r)^{-1}\right)\dd r\right]\quad(w>a).
\]
Writing \(\delta_{(a,a)}\) for the aligned atom, the limiting joint
measure on the initial right half is this atom with weight \(A/8\),
plus the measure on \(y=x>a\) with density in \(x\)
\[
 f_b(x)=\frac{A}{2\pi b}\frac{H'(x)}{(1+(H(x)/b)^2)^2}.
\]
Each sector has limiting mass \(A/8\). For \(\zeta>a\) the fraction of
the total \(A/4\) active at the current probe is therefore
\begin{equation}\label{an04rpsv1:fraction}
 \frac12+\frac2\pi\int_0^{H(\zeta)/b}\frac{\dd t}{(1+t^2)^2}.
\end{equation}
The atom is strictly active on the specified offset compact; the Q1
boundary carries no mass. Q2 and Q3 ancestry are excluded. This is
an entry comparison. No strong mass, phase or profile for the trained
model is identified with it.

\section{Rank transport, cutoff and moments}\label{an04rpsv1:rank}
Use normalized mass rank \(\sigma\in[0,1]\) with uniform rank law and
strong mass \(M(t)\). All labels have reaction \(1-M\), so normalized
rank weights remain fixed. Let \(Y=\int y\,d\mu,\ X_m=\int x\,d\mu,\
K_s=\int K(\rho x)d\mu,\ S(x)=\int F(\rho x,\rho\widetilde x)d\mu\),
where \(F(t,r)=\E[(Z+r)_+1_{Z+t>0}]\) for a standard normal \(Z\).
Use the actual current measure for every donor integral. Define
\[
 2U=-(1-M)x+\rho\phi(\rho x)-\rho S(x)+\lambda y\Phi(\rho x),
 \qquad 2V=-(1-M)y-Y+\rho K(\rho x).
\]
For prescribed compact nondecreasing entry coordinates the full
self-consistent solution preserves co-order: at each time the pair
difference solves the segment-averaged receiver Jacobian equation
with the same donors. Both off-diagonal entries are averages of
\(\lambda\Phi(\rho x)/2>0\). The integrating-factor argument of CFS
spread applies to every ordered pair at once. Bounded initial support
propagates on finite intervals by linear growth, as in the estimates
below. A plateau of identical pairs stays a plateau by uniqueness.
The active set is an initial segment of rank, of length
\(\sigma_*=\mu\{x<\zeta\}/M\). It equals \([0,\sigma_*)\) modulo
endpoints (take a right-continuous quantile convention on \([0,1)\)).
A plateau at equality is excluded in its entirety. This assertion
does not differentiate a moving rank boundary.

For cutoff choose \(R>\max(a,\zeta)\) large enough that
\(H(x)\ge x/2\) for \(x\ge R\), and set \(h_R=H(R)\).
Deleting the tail removes mass at most
\(Ab^3/(6\pi h_R^3)\) and first absolute input moment at most
\(Ab^3/(2\pi h_R^2)\), by \(x\le2H(x)\) and integration in
\(h=H(x)\). The same estimate holds for \(y\). The associated donor
errors in \(K_s,S(x)\) are bounded by \(\phi(0)\) times the mass
charge plus \(\rho\) times the first moment charge, uniformly in the
receiver. No exposure is removed at these offsets. One can instead
clamp this tail to \((R,R)\) and preserve mass and the active fraction.
The first moment change is then bounded by the same first moment charge,
and the kernel changes by at most \(\rho\) times that charge since
\(|F_r|,|K'|\le1\). Multiply charges by \(4M_0/A\) if normalizing
the comparison shape to a prescribed leading mass \(M_0\).
These are entry costs. The clamped evolution is a new self-consistent
solution, not a trajectory comparison with its unclamped counterpart.

For any such leading state the following equations are exact:
\begin{align}
 M'&=M(1-M),\qquad
 2Y'=(1-2M)Y+\rho K_s,\notag\\
 2X_m'&=(1-M)X_m+\int [\rho\phi(\rho x)-\rho S(x)+
                      \lambda y\Phi(\rho x)]\,d\mu,\notag\\
 K_s'&=(1-M)K_s+\rho\int\Phi(\rho x)U\,d\mu,\notag\\
 A_\zeta'&=(1-M)A_\zeta-\int_{x<\zeta}U\,d\mu,\quad
 A_\zeta=\int(\zeta-x)_+d\mu .
 \label{an04rpsv1:moments}
\end{align}
Apply transport with reaction to each integrand. The last line holds
almost everywhere in time: an absolutely continuous coordinate has
zero velocity almost everywhere on its level set, and finite-interval
bounded speeds allow integration over rank. At a gate atom the right
derivative has the additional term \(\int_{x=\zeta}(-U)_+d\mu\).
Exposure is distinct from active mass. For example, at \(M=1\) a
nonzero \(K_s\) forces the signed mean \(Y\); it cannot simply be set
to its coherent value during escape.

\section{A self-consistent escape region and its range}\label{an04rpsv1:escape}
Reset leading time to zero at a prescribed resident-phase entry.
Here is a sufficient test using only initial moments and the gate.
It is constructive on a nonempty range of the critical profiles.
\begin{theorem}\label{an04rpsv1:cone}
Let \(0<M_0\le1,\ N=0,\ \zeta>a\). For compact entry define
\[
 g=\rho K(\rho\zeta)-\max\{\rho K_s(0),Y(0)\}
                          -(1-M_0)(\zeta+1).
\]
When \(g>0\), let \(D_0=1+\int(|x|+|y|)d\mu_0\),
\(C=10D_0e^7\), and set
\begin{align*}
 k&=\frac{g}{8(\zeta+1)},&
 h&=\min\{1,g/(4C),\,k^{-1}\log((\zeta+1)/\zeta)\},\\
 \delta&=\min\{g/(4\lambda),(\zeta-a)/2,\,
                         \zeta(1-e^{-kh/2})/2\}.
\end{align*}
For each \(\xi\in[\zeta-\delta,\zeta]\) the region
\[
 0\le t\le h,\qquad x,y\ge w(t)=\xi e^{kt}
\]
is forward-invariant up to its stated terminal time. An aligned label
initially at \(\xi\in[\zeta-\delta,\zeta)\) leaves the active set by
\(h/2\) and stays inactive through \(h\), with both coordinates bounded
below by \(\xi e^{kt}\). For the clamped, mass-normalized critical
profile the fraction of initial mass so certified is at least
\begin{equation}\label{an04rpsv1:moved}
 \frac4A\int_{\zeta-\delta}^{\zeta} f_b(u)\,du>0.
\end{equation}
The barrier reaches radius \(\xi e^{kh}>\zeta\); for entry chart radius
\(R_0\) the entire proof takes place in
\([-(1+R_0)e^2,(1+R_0)e^2]^2\).
\end{theorem}
\begin{proof}
First obtain a donor budget without assuming any label stays below the
probe. Let \(D(t)=1+\int(|x|+|y|)d\mu\).
From \(0\le S(x)\le K_s\le M\phi(0)+\rho\int|x|d\mu\), \(|Y|\le D-1\)
and the field equations,
\[
 |U|+|V|\le3(1+|x|+|y|+D-1).
\]
The radial reaction is at most one, hence \(D^+D\le7D\).
This charges the escaping fraction with its own growing coordinates.
On \(t\le1\), \eqref{an04rpsv1:moments} yields
\(|Y'|,\rho|K_s'|\le C\) using \(\int|U|d\mu\le6D\).
Independently the support radius satisfies \(D^+R\le2(1+R)\)
by applying the same field formulas with \(K_s\le M(\phi(0)+\rho R)\).
This proves finite-interval continuation and the stated chart budget.

Now \(0<w\le\zeta+1\), \(M(t)\ge M_0\).
At the input face \(x=w,\ y\ge w\) one has
\[
 2U\ge -(1-M)w+\rho K(\rho w)-\rho K_s(t).
\]
Here both \(\rho\phi(\rho w)\) and \(\lambda w\Phi(\rho w)\)
are retained. At the output face \(y=w,\ x\ge w\),
\[
 2V\ge -(1-M)w+\rho K(\rho w)-Y(t).
\]
At either face the bound is at least \(g-\lambda\delta-Ch\ge g/2\),
using \((\rho K(\rho w))'\) with respect to \(w\) at most \(\lambda\).
Both velocities exceed \(w'=kw\le g/8\), including the corner.
The entry face starts at the prescribed label, and the terminal face
ends the estimate. Strict first-exit comparison proves the claim.
Its mass is positive by the CFS density; common reaction preserves
the fraction. This counts labels newly leaving the gate, independently
of the mass already above it.
\end{proof}

\paragraph{Parameters for which the test closes.}
For example take the normalized, clamped critical shape at \(M_0=1\)
with fixed \(R>\zeta\). Let \(b\downarrow0\) in this family of
\emph{entry states} before fixing a positive \(b\).
The Q1 rank coordinate can be written \(H^{-1}(b\tan\beta)\),
with normalized weight proportional to \(\cos^2\beta\) independent
of \(b\); the Q4 atom is fixed. Bounded convergence after clamping
gives \(Y(0)\to a\) and \(\rho K_s(0)\to\rho K(\rho a)=a\).
Thus \(g\to\rho K(\rho\zeta)-a>0\).
There is also a simple sufficient parameter interval with no moment
integrals to evaluate. Put \(r_{\rm mid}=[a+\rho K(\rho\zeta)]/2\);
then \(a<r_{\rm mid}<\rho K(\rho\zeta)<\zeta\). At \(M_0=1\)
the mass above \(r_{\rm mid}\) is at most
\(2b^3/[3\pi H(r_{\rm mid})^3]\) by integration of \(g_b\).
For aligned clamped data, \(\rho K(\rho u_R)\le u_R\) for \(u_R\ge a\).
Hence the test succeeds for every
\[
 0<b<H(r_{\rm mid})
 \left(\frac{3\pi[\rho K(\rho\zeta)-r_{\rm mid}]}{2R}\right)^{1/3}
 \qquad \text{at } M_0=1:
\]
indeed \(Y(0)\le r_{\rm mid}+2Rb^3/[3\pi H(r_{\rm mid})^3]\).
An explicit test for any fixed \(b\) uses the integrals of
\(u_R=\min(u,R)\) and \(\rho K(\rho u_R)\) against
\((4M_0/A)[(A/8)\delta_a+f_b(u)du]\) in the definition of \(g\).
By continuity the same strict test holds at sufficiently close
subunit masses in this range. This is a leading escape result at
fixed parameters, without a coherent-donor approximation.
If a later resident state was reached dynamically, its own current
moments and available ordered pairs are required in this test.

\paragraph{Obstruction for unrestricted critical parameters.}
Even at \(M=1\), the exact input-face test at \(x=\zeta,\ y=\zeta\)
for an aligned initial state is
\begin{equation}\label{an04rpsv1:face}
 2U(\zeta,\zeta)=\rho[K(\rho\zeta)-S(\zeta)] .
\end{equation}
Its sign is not fixed. For the same normalization and cutoff construction,
as \(b\to\infty\) the Q1 half converges to the clamped point \(R\).
Thus \(S(\zeta)\to [K(\rho a)+\phi(\rho\zeta)+
\rho R\Phi(\rho\zeta)]/2\), since \(a<\zeta<R\).
Choose a finite \(R>\zeta\) large enough that this exceeds
\(K(\rho\zeta)\), then fix sufficiently large finite \(b\).
One can choose such an \(R\) first even without a tail-smallness
requirement: clamping is a prescribed compact leading model. The
negative-face example also holds with an arbitrarily small removed
tail at fixed large \(b\): increasing \(R\) only increases the
initial \(S(\zeta)\), because \(F_r\ge0\). The explicit cutoff
charges then apply. At this input face the outward test fails even
with zero prescribed barrier speed. The aligned gate has positive
rank density in a neighborhood on either side. By continuity, nearby
initially inactive labels move into the active segment at short time.
This calculation precludes an unrestricted outward-gate argument.
It does not decide eventual escape or capture for this latter family.
CFS's output-gap bound separately precludes collapse of separated
output quantiles into one coherent point. It does not place the input
quantiles relative to a specified \(\zeta\).

\paragraph{Active-rank monotonicity.}
At a smooth gate with positive \(x_\sigma\),
\(\sigma_*'=-U(\zeta,y(\sigma_*))/x_\sigma\).
In particular the clamped shapes above have a regular gate initially.
Equation \eqref{an04rpsv1:face} gives negative rank derivative for
the small-\(b\) example and positive derivative for the large-\(b\)
example. A parameter-uniform monotonic direction is impossible.
Within any particular field a proof of perpetual nonincrease would
need the input-face inequality
\[
 \rho\phi(\rho\zeta)-\rho S(\zeta)+\lambda y(\sigma_*)\Phi(\rho\zeta)
 \ge(1-M)\zeta
\]
at the moving gate, with separate handling of equality atoms.
Rank order does not supply it. The examples compare signs computed
within their respective fields. They do not prove temporal
nonmonotonicity along either single trajectory.

\section{Handoff}\label{an04rpsv1:handoff}
S7 has a right-half comparison entry measure and an \(N=0\) compact
leading evolution in rank coordinates, with evolving \(M,Y,X_m,K_s\)
and \(A_\zeta\). For the range of Theorem~\ref{an04rpsv1:cone} there
is a certified newly inactive rank band and a lower bound through \(h\).
For arbitrary prescribed critical parameters, the later split and even
the sign of motion at the gate must still be determined from the field.
There is no identified universal end time for the resident phase.

S7b needs a matched switch to weak growth, the actual seed and joint
weak law (the candidate scale is \(N_0\asymp S_0^{1-\lambda}\)),
strong and weak donor feedback and radial weights through passage,
and the current gate, exact signed lag and rate-weighted producer
bounds at the probe scale. The transfer obligations remain the
trained entry clock, \(o(q)\)-scale angular resolution or equivalent
weighted-profile accuracy, radial mass and other-sector/tail budgets,
and comparison through strong saturation in the full trained field.
The critical comparison ratio \((S_0/A)^\omega/q=b\) is fixed and
the deep all-right-half overlap has the wrong power here.
\end{document}
